\documentclass[12pt,a4paper]{article}
\usepackage[margin=25mm, headheight=15pt, headsep=10pt]{geometry}
\usepackage{fontspec}
\usepackage[table]{xcolor} % Note: table option is required for rowcolors
\usepackage{titlesec}
\usepackage{tocloft}
\usepackage{fancyhdr}
\usepackage{booktabs}
\usepackage{tcolorbox}
\tcbuselibrary{skins, breakable}
\usepackage[colorlinks=true, linkcolor=emerald, urlcolor=emerald, bookmarks=true]{hyperref}

\definecolor{emerald}{RGB}{0,102,51}
\definecolor{darkred}{RGB}{139,0,0}
\definecolor{lightgray}{RGB}{245,245,245}

\usepackage[nil]{babel}
\babelprovide[import, main]{bengali}
\babelprovide[import]{arabic}
\babelfont{rm}[Renderer=HarfBuzz,Scale=1.0]{Noto Serif Bengali}
\babelfont{sf}[Renderer=HarfBuzz]{Noto Sans Bengali}
\babelfont{tt}[Renderer=HarfBuzz]{Noto Sans Bengali}
\babelfont[arabic]{rm}[Renderer=HarfBuzz,Scale=1.1]{Noto Naskh Arabic}

% Section styling
\titleformat{\section}{\Large\bfseries\color{emerald}}{\thesection.}{0.5em}{}
\titleformat{\subsection}{\large\bfseries\color{emerald!85!black}}{\thesubsection.}{0.5em}{}
\titleformat{\subsubsection}{\normalsize\bfseries\color{emerald!70!black}}{\thesubsubsection.}{0.5em}{}
\titlespacing*{\section}{0pt}{18pt}{8pt}

% Fancy Header/Footer setup
\pagestyle{fancy}
\fancyhf{} % clear all header and footer fields
\fancyhead[L]{\color{emerald}\small\textbf{\nouppercase{\leftmark}}}
\fancyfoot[L]{\scriptsize\color{black!50}{\lat{AI}}-সংকলিত, আলেম-যাচাইকৃত নয়}
\fancyfoot[C]{\color{emerald!70!black}\thepage}
\renewcommand{\headrulewidth}{0.4pt}
\renewcommand{\headrule}{\hbox to\headwidth{\color{emerald}\leaders\hrule height \headrulewidth\hfill}}
\renewcommand{\footrulewidth}{0pt}

% Custom boxes
% 1. Ayat/Hadith Box
\newtcolorbox{ayatbox}[1][]{
    enhanced, breakable, 
    colback=emerald!5, 
    colframe=emerald, 
    boxrule=0.5pt, 
    arc=3pt, 
    left=10pt, right=10pt, top=10pt, bottom=10pt,
    #1
}

% 2. Quote/Translation Box
\newtcolorbox{quotebox}[1][]{
    enhanced, breakable,
    frame hidden,
    colback=lightgray,
    borderline west={3pt}{0pt}{emerald},
    left=15pt, right=10pt, top=10pt, bottom=10pt,
    sharp corners,
    #1
}

% 3. AI-generated notice box
\newtcolorbox{warnbox}[1][]{
    enhanced, breakable,
    colback=red!4,
    colframe=darkred,
    boxrule=0.8pt,
    arc=3pt,
    left=10pt, right=10pt, top=10pt, bottom=10pt,
    #1
}

% Commands
\newcommand{\arab}[1]{\foreignlanguage{arabic}{#1}}
\newcommand{\kib}[1]{\textcolor{emerald}{\textbf{#1}}}

% Latin (English) fallback: Noto Serif Bengali has NO Latin glyphs
\newfontfamily\latfont[Renderer=HarfBuzz]{Noto Serif}
\newcommand{\lat}[1]{{\latfont #1}}

\setlength{\parskip}{5pt}
\setlength{\parindent}{0pt}

% Fix for missing bullet in Noto Serif Bengali
\renewcommand{\labelitemi}{$\bullet$}



\begin{document}
\begin{center}{\Large\bfseries\color{emerald} ক্ষমা করা পাপ — \lat{02}-আমর-ইবনে-আস-মৃত্যুশয্যায়-আশা}\end{center}
\vspace{1em}
\section*{আমর ইবনে আস (রাঃ) — মৃত্যুশয্যায় অতীতের জন্য আশা}

\subsection*{সূত্র কার্ড}

\begin{table}[h]\centering\small
\begin{tabular}{ll}
বিষয় & বিবরণ \\
\hline
\textbf{ঘটনা} & আমর ইবনে আস (রাঃ) মৃত্যুশয্যায় কাঁদেন; ইসলাম ও হিজরত পূর্বের পাপ মুছে দেওয়ার নবীজির (সা.) ঘোষণা স্মরণ করেন \\
\textbf{বর্ণনাকারী} & ইবনে শামাসা মাহরি (রহ.) — আমর ইবনে আসের (রাঃ) মৃত্যুর সময় উপস্থিত \\
\textbf{গ্রন্থ ও নম্বর} & সহীহ মুসলিম ১২১ \\
\textbf{অধ্যায়} & কিতাবুল ঈমান — \lat{“}ইসলাম যা আগে ছিল তা ধ্বংস করে দেয়, হিজরত ও হজও\lat{”} \\
\textbf{সনদের মান} & \textbf{সহীহ} \\
\textbf{স্থান ও সময়} & মিসর — আমর ইবনে আসের (রাঃ) মৃত্যুশয্যা (হিজরি ৪৩ সালের দিকে) \\
\hline
\end{tabular}
\end{table}

\subsection*{সংক্ষিপ্ত পটভূমি}

আমর ইবনে আস (রাঃ) ইসলামের আগে কুরাইশের একজন শক্তিশালী ও প্রভাবশালী ব্যক্তি ছিলেন — ইসলাম ও নবীজির (সা.) বিরুদ্ধে সক্রিয় ছিলেন। ইসলাম গ্রহণের পর তিনি হয়ে ওঠেন নবীজির (সা.) প্রিয় সাহাবি, বিজয়ী সেনাপতি (মিসর বিজয়) ও বড় আলেম। মৃত্যুর সময় তার নিজের অতীতের কথা মনে পড়ে — এবং তিনি নবীজির (সা.) সেই ঘোষণার কথা স্মরণ করেন, যা অতীতের পাপের জন্য আশার দরজা খুলে দেয়।

\subsection*{ঘটনার পূর্ণ বিবরণ}

ইবনে শামাসা মাহরি (রহ.) বলেন: আমরা আমর ইবনে আসের (রাঃ) কাছে ছিলাম — তিনি মৃত্যুশয্যায়। তিনি অনেকক্ষণ কাঁদলেন এবং মুখ দেয়ালের দিকে ঘুরিয়ে নিলেন।

তাঁর ছেলে (আবদুল্লাহ) বললেন: \lat{“}বাবা! রাসূলুল্লাহ (সা.) কি আপনাকে সুসংবাদ দেননি? রাসূলুল্লাহ (সা.) কি আপনাকে সুসংবাদ দেননি?\lat{”}

তিনি মুখ ঘুরিয়ে বললেন: \lat{“}হ্যাঁ — সর্বোত্তম জিনিস যা আমরা ভরসা করি তা হলো — আল্লাহ ছাড়া কোনো ইলাহ নেই এবং মুহাম্মাদ আল্লাহর রাসূল এই সাক্ষ্য। আমি তিনটি অবস্থা অতিক্রম করেছি:

\begin{enumerate}
\item \textbf{প্রথম অবস্থা:} রাসূলুল্লাহকে (সা.) আমি যতটা ঘৃণা করতাম, পৃথিবীতে আর কাউকে ততটা ঘৃণা করতাম না; আমার একটাই ইচ্ছা ছিল — তাকে পরাস্ত করা ও হত্যা করা। এই অবস্থায় মারা গেলে আমি নিশ্চিত জাহান্নামি হতাম।
\end{enumerate}
\begin{enumerate}
\item \textbf{দ্বিতীয় অবস্থা:} আল্লাহ আমার অন্তরে ইসলামের ভালোবাসা দান করলেন। আমি নবীজির (সা.) কাছে এসে বললাম: আপনার ডান হাত বাড়ান, আমি আপনার কাছে বাইয়াত করি। তিনি হাত বাড়ালেন — আমি হাত গুটিয়ে নিলাম। তিনি বললেন: "কী হয়েছে?" আমি বললাম: আমি শর্ত রাখতে চাই। তিনি বললেন: "কী শর্ত?" আমি বললাম: আমি যেন ক্ষমা পাই। তখন তিনি বললেন: \textbf{"তুমি কি জানো না যে, ইসলাম যা আগে ছিল তা মুছে দেয়? হিজরত যা আগে ছিল তা মুছে দেয়? আর হজ যা আগে ছিল তা মুছে দেয়?"}
\end{enumerate}
\begin{enumerate}
\item \textbf{তৃতীয় অবস্থা:} এরপর আমার কাছে রাসূলুল্লাহ (সা.)-এর চেয়ে প্রিয় আর কেউ ছিল না; তাঁর চেয়ে মহান চোখে আর কেউ ছিল না। আমি তাঁর মুখের দিকে তাকাতেই পারতাম না — তাঁর দীপ্তিময়তার কারণে। আমাকে জিজ্ঞেস করলে তাঁর চেহারা বর্ণনা করতে পারব না — কারণ আমি কখনো পুরোপুরি তাঁর দিকে তাকাইনি। এই অবস্থায় মারা গেলে আমার জান্নাতি হওয়ার আশা থাকত।
\end{enumerate}
এরপর আমাদের থেকে কিছু বিষয় ঘটে গেছে (ফিতনার ঘটনা) — এখন আমি জানি না আমার জন্য কী আছে। যখন আমি মারা যাব, আমার সাথে কোনো বিলাপকারিণী বা আগুন (শোকের রীতি) যেন না থাকে। যখন আমাকে কবর দেবে, মাটি ভালোভাবে চাপা দিয়ে কবরের পাশে দাঁড়িয়ে থাকো — উট জবাই ও গোশত বণ্টনের সময় পরিমাণ — যেন আমি তোমাদের সাহচর্যে থাকতে পারি এবং আল্লাহর ফেরেশতাদের জবাব দিতে প্রস্তুত হতে পারি।\lat{”}

(সহীহ মুসলিম ১২১)

\subsection*{শব্দের অর্থ}

\begin{table}[h]\centering\small
\begin{tabular}{ll}
শব্দ & অর্থ \\
\hline
\textbf{ইয়াহদিমু মা কাবলাহু (\arab{يهدم} \arab{ما} \arab{قبله})} & যা আগে ছিল তা ধ্বংস/মুছে দেয় \\
\textbf{আল-বাইয়াত (\arab{البيعة})} & আনুগত্যের শপথ \\
\textbf{সিয়াকাতুল মাওত} & মৃত্যুশয্যা — প্রাণ কণ্ঠাগত অবস্থা \\
\hline
\end{tabular}
\end{table}

\subsection*{সংশ্লিষ্ট দলিল}

\begin{itemize}
\item \textbf{আয়াত (২৫:৭০):} তাওবার পর পাপ নেকিতে বদল — তাওবার আগের পাপের দরজা খোলা।
\item \textbf{হাদিস (বুখারি ২৪৪১; মুসলিম ২৭৬৮):} নাজওয়া — কিয়ামতে মুমিনের পাপ একান্তে স্মরণ করিয়ে ক্ষমা।
\item \textbf{হাদিস (বুখারি ৬৫৩৬):} মুমিনের সহজ হিসাব = আরয।
\end{itemize}
\subsection*{শিক্ষা}

\begin{enumerate}
\item \textbf{অতীত যত অন্ধকারই হোক — আল্লাহর দিকে ফেরার দরজা খোলা।} আমর ইবনে আস (রাঃ) ইসলামের ঘোর শত্রু ছিলেন — তবু ইসলাম গ্রহণের পর তার অতীত মুছে গেল।
\item \textbf{নবীজির (সা.) ঘোষণা — ইসলাম/হিজরত/হজ অতীত মুছে দেয়} — নতুন মুসলিম ও তাওবাকারীর জন্য আশার বড় দলিল।
\item আমর ইবনে আস (রাঃ) মৃত্যুর সময়ও নিজের হিসাব নিয়ে শঙ্কিত ছিলেন — \textbf{সাহাবিরাও হিসাবের ভয় করতেন; কিন্তু সেই ভয় হতাশার নয়, আল্লাহর দিকে ফেরার।} তাওবাকারীর ভয় আল্লাহর রহমতের আশা থেকে পৃথক নয়।
\end{enumerate}

\end{document}
